\documentclass[10pt]{article}
\usepackage[margin=1.05in]{geometry}
\usepackage{amsmath,amssymb,amsthm}
\usepackage[T1]{fontenc}
\allowdisplaybreaks

\newtheorem{theorem}{Theorem}[section]
\newtheorem{proposition}[theorem]{Proposition}
\newtheorem{lemma}[theorem]{Lemma}
\newtheorem{corollary}[theorem]{Corollary}
\newtheorem{definition}[theorem]{Definition}
\theoremstyle{remark}
\newtheorem{remark}[theorem]{Remark}

\newcommand{\A}{\mathcal{A}}
\newcommand{\R}{\mathbb{R}}
\newcommand{\Om}{\Omega}
\newcommand{\Sig}{\Sigma}
\newcommand{\Gam}{\Gamma}
\newcommand{\Pmax}{P_{\max}}
\newcommand{\Pmin}{P_{\min}}
\newcommand{\dE}{d_{\mathrm{E}}}
\newcommand{\dP}{d_{\Om}}
\newcommand{\dA}{d_{\A}}
\newcommand{\norminf}[1]{\|#1\|_{L^\infty}}
\newcommand{\eps}{\varepsilon}

\title{Forced Mechanics of the Mixed-Boundary Torsion Problem:\\
\large Exact Laws, Sharp Star-Shaped Bounds, and the Closure of the Distance
Conjecture}
\author{Echo-S Research}
\date{July 18, 2026}

\begin{document}
\maketitle

\begin{abstract}
This paper records, with complete proofs, the forced mechanics of the
mixed-boundary torsion problem $-\nabla\!\cdot\!(\A\nabla u)=q$ with $u=0$ on a
supply portion $\Sig\subset\partial\Om$ and the natural (conormal Neumann)
condition on the remainder $\Gam$; coefficients are measurable, symmetric,
and uniformly elliptic, $\Pmin I\le\A\le\Pmax I$ a.e., with every constant
below depending on $\Pmax$ only. The results: (1) the exact one-dimensional
law $u(x)=q\int_0^x (L-s)/A(s)\,ds$, with the sharp far-end bound
$qL^2/(2\Pmax)$ and equality exactly for $A\equiv\Pmax$; (2) the ball-comparison
mean bound: for every ball $B_r\subseteq\Om$,
$\norminf{u}\ge q\,r^2/\bigl(n(n+2)\Pmax\bigr)$ --- constant $15$ in dimension
three --- via same-operator comparison and torsional-rigidity anti-monotonicity;
(3) the pointwise full-boundary bound
$u(x)\ge q\,\mathrm{dist}(x,\partial\Om)^2/(2n\Pmax)$ for constant tensors;
(4) the sharp star-shaped theorem: if $\Gam$ is star-shaped about $x_0$, then
$u(x_0)\ge (q/2n)\inf_{\Sig}(z-x_0)^{\!\top}\A^{-1}(z-x_0)\ge
q\,\dE(x_0,\Sig)^2/(2n\Pmax)$, ratio-free, with equality on balls, central
reflection sectors, and ellipsoids; (5) the closure of the distance conjecture
P2: an exactly solvable family of tapered horns yields
$\Pmax\norminf{u}/(qR^2)\le
(1-\delta^2+k\eps^2/2)/\bigl((2+k(n-1))(1-\delta)^2\bigr)$, so all three metric
variants of the conjectured bound $qR^2/(2n\Pmax)$, $R=\sup_x d(x,\Sig)$, are
false, the infimum of the coefficient over admissible data is $0$, and an
explicit Lipschitz instance attains the exact rational value bound
$423/2890<1/6$ in dimension three; and (6) the dichotomy within the horn family: the star
condition degenerates exactly at the straight cone $k=2$, where the
coefficient equals $1/(2n)$ in every dimension --- a sharp \emph{sufficient}
hypothesis whose removal beyond the cone permits arbitrarily severe failure
(necessity over all geometries is not asserted). Distance to the full
boundary always controls $u$ from below; distance to the supply portion does
so under the star condition; on the horns beyond it, the governing quantity is
the volume-graded resistance, for which they provide the exact identity. Only
forced material appears; every displayed identity and rational inequality has
additionally been machine-verified in exact arithmetic.
\end{abstract}

\section{Setting and conventions}\label{sec:setting}

\paragraph{Standing conventions.}
Throughout, $n\ge 2$ unless stated otherwise, $q>0$ is constant, and $\A$
denotes a measurable symmetric coefficient tensor field which is uniformly
elliptic: $\Pmin I\le\A\le\Pmax I$ a.e.\ for fixed constants
$0<\Pmin\le\Pmax$. The lower bound enters only through well-posedness
(coercivity); no constant in any estimate below depends on it. We freely use classical facts --- Lax--Milgram, Rellich
compactness, trace theory on Lipschitz boundaries, Stampacchia's truncation
identities, and interior elliptic regularity for constant-coefficient
operators --- and prove everything else. ``Ratio-free'' means that a constant
depends on $\Pmax$ only, never on the ellipticity ratio $\Pmax/\Pmin$, the
volume, or regularity moduli.

\begin{definition}[Mixed torsion problem, trace formulation]\label{def:weak}
Let $\Om\subset\R^n$ be a bounded connected Lipschitz domain and let
$\partial\Om=\overline\Sig\cup\overline\Gam$ with $\Sig$ relatively open of
positive surface measure, and let $\A$ be measurable, symmetric, and uniformly
elliptic ($\Pmin I\le\A\le\Pmax I$ a.e., $\Pmin>0$). Set
\[
V \;=\; \{\varphi\in H^1(\Om)\;:\;\operatorname{tr}\varphi=0
\ \text{$\sigma$-a.e.\ on }\Sig\}.
\]
The \emph{torsion function} $u\in V$ is the unique solution of
\begin{equation}\label{eq:weak}
B[u,\varphi]\;:=\;\int_\Om \nabla\varphi\cdot\A\nabla u\,dz
\;=\;\int_\Om q\,\varphi\,dz\qquad\text{for all }\varphi\in V,
\end{equation}
which exists by the standard Poincar\'e inequality on $V$, coercivity from
uniform ellipticity, and Lax--Milgram.
Formulation \eqref{eq:weak} encodes $u=0$ on $\Sig$ and the natural condition
$\A\nabla u\cdot\nu=0$ on $\Gam$.
\end{definition}

\begin{lemma}[Positivity and nontriviality]\label{lem:pos}
$u\ge 0$ a.e., and $u\not\equiv 0$.
\end{lemma}

\begin{proof}
$u^-:=\max(-u,0)\in V$ (its trace vanishes on $\Sig$). Then
$B[u,u^-]=-\int_{\{u<0\}}\nabla u\cdot\A\nabla u\le 0$ while
$B[u,u^-]=\int q\,u^-\ge 0$; hence both vanish, $\nabla u^-\equiv 0$, and the
constant $u^-$ has zero trace on $\Sig$, so $u^-\equiv 0$. If $u\equiv 0$ then
\eqref{eq:weak} forces $\int q\varphi=0$ for all $\varphi\in V$, impossible.
\end{proof}

\paragraph{Metrics and the distance conjecture.}
Write $\dE(z,\Sig)=\inf_{w\in\Sig}|z-w|$ for the Euclidean distance,
$\dP(z,\Sig)$ for the infimal length of rectifiable paths inside $\Om$ from $z$
to $\Sig$, and, for locally Lipschitz $f$ on $\Om$,
\begin{equation}\label{eq:intrinsic}
\dA(z,\Sig)\;=\;\sup\Bigl\{f(z)\;:\;\nabla f\cdot\A\nabla f\le 1\ \text{a.e.},
\ \ \limsup_{w\to\Sig} f(w)\le 0\Bigr\}
\end{equation}
for the operator-intrinsic distance. With
$R_\bullet=\sup_{z\in\Om}d_\bullet(z,\Sig)$, the \emph{distance conjecture} P2
asserted, for all admissible data,
\begin{equation}\label{eq:P2}
\text{(P2-E)}\ \ \norminf{u}\ \ge\ \frac{q\,R_{\mathrm E}^2}{2n\,\Pmax},\qquad
\text{(P2-G)}\ \ \norminf{u}\ \ge\ \frac{q\,R_{\Om}^2}{2n\,\Pmax},\qquad
\text{(P2-A)}\ \ \norminf{u}\ \ge\ \frac{q\,R_{\A}^2}{2n}.
\end{equation}
Since $\dE\le\dP\le\sqrt{\Pmax}\,\dA$, the variants are nested:
P2-A $\Rightarrow$ P2-G $\Rightarrow$ P2-E. Sections
\ref{sec:oned}--\ref{sec:star} establish what \emph{does} hold; Sections
\ref{sec:horn}--\ref{sec:dichotomy} close \eqref{eq:P2} and exhibit, within an
explicit family, the exact transition between the star-controlled regime and
unbounded failure.

\begin{lemma}[Metric collapse for constant scalar coefficients]\label{lem:metric}
If $\A=P\,I$ with constant $P>0$, then $\dA(z,\Sig)=\dP(z,\Sig)/\sqrt P$ for
every $z\in\Om$.
\end{lemma}

\begin{proof}
``$\ge$'': $f_0:=\dP(\cdot,\Sig)/\sqrt P$ is admissible in
\eqref{eq:intrinsic}: small balls in $\Om$ are convex, so $\dP$ is locally
$1$-Lipschitz, hence $|\nabla f_0|\le P^{-1/2}$ a.e., i.e.\
$\nabla f_0\cdot\A\nabla f_0\le 1$; and $f_0\to 0$ at $\Sig$.
``$\le$'': any admissible $f$ has $|\nabla f|\le P^{-1/2}$ a.e., hence is
$P^{-1/2}$-Lipschitz along every rectifiable path in $\Om$; integrating along
paths from $z$ that approach $\Sig$ and using $\limsup_{\Sig}f\le 0$ gives
$f(z)\le \dP(z,\Sig)/\sqrt P$.
\end{proof}

\section{The one-dimensional law}\label{sec:oned}

\begin{theorem}[Exact 1D law]\label{thm:oned}
Let $0<\Pmin\le A(x)\le\Pmax$ be measurable on $(0,L)$ and $q>0$. The problem
\[
-\bigl(A u'\bigr)'=q\ \text{ on }(0,L),\qquad u(0)=0,\qquad (Au')(L)=0,
\]
has the unique solution (with $Au'$ absolutely continuous)
\[
u(x)\;=\;q\int_0^x \frac{L-s}{A(s)}\,ds .
\]
Consequently $u$ is nondecreasing,
\[
u(x)\;\ge\;\frac{q\,x^2}{2\Pmax}\quad\text{for every }x\in[0,L],
\qquad
\norminf{u}\;=\;u(L)\;=\;q\int_0^L\frac{L-s}{A(s)}\,ds\;\ge\;\frac{qL^2}{2\Pmax},
\]
and equality holds in the far-end bound if and only if $A=\Pmax$ a.e.\ on
$(0,L)$.
\end{theorem}

\begin{proof}
\emph{Existence.} With $u$ as displayed, $A(x)u'(x)=q(L-x)$ by construction;
this is absolutely continuous, vanishes at $x=L$, and $-(Au')'=q$ a.e.;
$u(0)=0$.
\emph{Uniqueness.} A difference $v$ of two solutions satisfies $(Av')'=0$, so
$Av'\equiv c$; the far-end condition gives $c=0$, hence $v'=0$ a.e.\ and
$v\equiv v(0)=0$.
\emph{Monotonicity and bounds.} $u'=q(L-x)/A\ge0$. Since $A\le\Pmax$ and
$L-s\ge x-s$ on $[0,x]$,
$u(x)\ge \tfrac{q}{\Pmax}\int_0^x(x-s)\,ds=qx^2/(2\Pmax)$; at $x=L$,
$u(L)\ge \tfrac{q}{\Pmax}\int_0^L(L-s)\,ds=qL^2/(2\Pmax)$, with equality iff
$\int_0^L(L-s)\bigl(\tfrac1A-\tfrac1{\Pmax}\bigr)ds=0$; the integrand is
nonnegative and $L-s>0$ a.e., forcing $A=\Pmax$ a.e.
\end{proof}

\begin{remark}\label{rem:onedconst}
Here $\Sig=\{0\}$ and $d(x,\Sig)=x$, so the pointwise bound reads
$u\ge q\,d(\cdot,\Sig)^2/(2\Pmax)$: the constant matches $2n$ at $n=1$; the
pointwise bound is strict for every $x<L$. The
sink in Theorem~\ref{thm:oned} is $q$ per unit \emph{length}; the
volume-graded 1D balance, with sink $q$ per unit \emph{volume} in a tube of
varying cross-section, is a different law and is the mechanism of
Section~\ref{sec:horn} (see Remark~\ref{rem:mechanism}).
\end{remark}

\section{Ball comparison and the mean bound}\label{sec:ball}

\begin{lemma}[Same-operator domain comparison]\label{lem:compare}
Let $\omega\subseteq\Om$ be open and let $w\in H^1_0(\omega)$ be the
full-Dirichlet torsion function of $\omega$ for the \emph{same} coefficient
$\A$ and sink $q$:
$\int_\omega\nabla\varphi\cdot\A\nabla w=\int_\omega q\varphi$ for all
$\varphi\in H^1_0(\omega)$. Then $u\ge w$ a.e.\ on $\omega$.
\end{lemma}

\begin{proof}
Let $w_j\in C_c^\infty(\omega)$ with $w_j\to w$ in $H^1_0(\omega)$. Since
$u\ge0$ (Lemma~\ref{lem:pos}), $(w_j-u)^+$ vanishes wherever $w_j=0$, so it is
an $H^1$ function supported in $\operatorname{supp}w_j\Subset\omega$; hence
$(w_j-u)^+\in H^1_0(\omega)$, and by truncation continuity
$\varphi:=(w-u)^+\in H^1_0(\omega)$. Its extension by zero to $\Om$ lies in
$V$. Testing \eqref{eq:weak} with the extension and the $\omega$-problem with
$\varphi$ and subtracting,
\[
0\;=\;\int_\omega\nabla\varphi\cdot\A\nabla(w-u)
\;=\;\int_{\{w>u\}}\nabla(w-u)\cdot\A\nabla(w-u)\;\ge\;0,
\]
so $\nabla\varphi\equiv0$ and $\varphi\in H^1_0(\omega)$ is constant, hence
$\varphi\equiv0$: $w\le u$ a.e.\ on $\omega$.
\end{proof}

\begin{lemma}[Torsional-rigidity anti-monotonicity]\label{lem:rigidity}
Let $\omega$ be bounded open and let $w,w'$ be the full-Dirichlet torsion
functions of $\omega$ for coefficients $\A\le\A'$ (a.e., as quadratic forms)
and the same $q$. Then $\int_\omega w\;\ge\;\int_\omega w'$.
\end{lemma}

\begin{proof}
For $\varphi\in H^1_0(\omega)$ set
$J_{\A}(\varphi)=2q\int\varphi-\int\nabla\varphi\cdot\A\nabla\varphi$. Since
$B_\A[w,\varphi]=q\int\varphi$,
\[
J_\A(\varphi)\;=\;2B_\A[w,\varphi]-B_\A[\varphi,\varphi]
\;=\;B_\A[w,w]-B_\A[\varphi-w,\varphi-w]
\;\le\;B_\A[w,w]\;=\;q\int_\omega w ,
\]
with equality at $\varphi=w$: thus $q\int w=\max J_\A$. Since
$J_\A(\varphi)\ge J_{\A'}(\varphi)$ pointwise in $\varphi$,
$q\int w=\max J_\A\ge\max J_{\A'}=q\int w'$.
\end{proof}

\begin{theorem}[Ball mean bound; the P1 constant $n(n+2)$]\label{thm:P1}
Let $u$ be as in Definition~\ref{def:weak}, with measurable, uniformly
elliptic $\A\le\Pmax I$. For every open ball $B_r(x_0)\subseteq\Om$,
\[
\frac{1}{|B_r|}\int_{B_r(x_0)}u\;\ge\;\frac{q\,r^2}{n(n+2)\,\Pmax},
\qquad\text{hence}\qquad
\norminf{u}\;\ge\;\frac{q\,r^2}{n(n+2)\,\Pmax}.
\]
In particular $\norminf{u}\ge q\,r_{\mathrm{in}}(\Om)^2/(n(n+2)\Pmax)$ with
$r_{\mathrm{in}}$ the inradius, which is attained:
$\mathrm{dist}(\cdot,\partial\Om)$ is continuous on the compact
$\overline\Om$, so the open ball of radius $r_{\mathrm{in}}$ about a maximizer
lies in $\Om$. For $n=3$ the constant is $15$.
\end{theorem}

\begin{proof}
Let $w_\A$ be the full-Dirichlet torsion of $B:=B_r(x_0)$ for $\A$, and
$w_{\Pmax}$ that for $\Pmax I$, namely
\[
w_{\Pmax}(z)\;=\;\frac{q\,(r^2-|z-x_0|^2)}{2n\Pmax}
\]
(direct verification: $-\Pmax\Delta w_{\Pmax}=q$, zero boundary values).
By Lemma~\ref{lem:compare}, $u\ge w_\A$ on $B$; by
Lemma~\ref{lem:rigidity} with $\A'=\Pmax I$,
$\int_B w_\A\ge\int_B w_{\Pmax}$. The radial mean of $r^2-|z-x_0|^2$ over $B$
is
\[
\frac{\int_0^r (r^2-\rho^2)\rho^{n-1}d\rho}{\int_0^r \rho^{n-1}d\rho}
=\frac{r^{n+2}\bigl(\tfrac1n-\tfrac1{n+2}\bigr)}{r^n/n}
=\frac{2r^2}{n+2},
\]
so $\frac1{|B|}\int_B u\ge\frac1{|B|}\int_B w_{\Pmax}
=\frac{q}{2n\Pmax}\cdot\frac{2r^2}{n+2}=\frac{qr^2}{n(n+2)\Pmax}$, and the
essential supremum dominates the mean.
\end{proof}

\begin{remark}\label{rem:P1scope}
Only $B\subseteq\Om$ is used: the ball may touch $\partial\Om$, including
$\Sig$, and the coefficient may be any measurable uniformly elliptic tensor. The bound is ratio-free
and mean-form, not pointwise; the pointwise sharp bound for \emph{constant}
tensors is Corollary~\ref{cor:fullbdry} below.
\end{remark}

\section{The star-shaped theorem}\label{sec:star}

Throughout this section $\Om$ is a bounded connected Lipschitz domain as in
Definition~\ref{def:weak}, and $\A$ is a \emph{constant} symmetric positive
definite tensor with $\A\le\Pmax I$; constancy makes uniform ellipticity
automatic, with $\Pmin=\lambda_{\min}(\A)$.

\begin{proposition}[Two identities]\label{prop:f}
Fix $x_0\in\Om$ and set $f(z)=(z-x_0)^{\!\top}\A^{-1}(z-x_0)$. Then
\[
\A\nabla f=2(z-x_0),\qquad
\nabla\!\cdot\!(\A\nabla f)=2n,\qquad
\A\nabla f\cdot\nu=2(z-x_0)\cdot\nu\ \ \text{on }\partial\Om,
\]
and $f(z)\ge|z-x_0|^2/\Pmax$, since in an orthonormal eigenbasis of $\A$ with
eigenvalues $\lambda_i\le\Pmax$,
$f-|z-x_0|^2/\Pmax=\sum_i w_i^2(\Pmax-\lambda_i)/(\lambda_i\Pmax)\ge0$.
\end{proposition}

\begin{proof}
$\nabla f=2\A^{-1}(z-x_0)$, hence $\A\nabla f=2(z-x_0)$, whose divergence is
$2n$; the conormal identity and the eigenbasis decomposition are immediate.
\end{proof}

\begin{theorem}[Star-shaped lower bound]\label{thm:B}
Let $u$ be the torsion function of Definition~\ref{def:weak} and let
$x_0\in\Om$ satisfy the star condition
\begin{equation}\label{eq:star}
(z-x_0)\cdot\nu(z)\;\ge\;0\qquad\text{for $\sigma$-a.e. }z\in\Gam .
\end{equation}
Then
\[
u(x_0)\;\ge\;\frac{q}{2n}\,\inf_{z\in\Sig}
(z-x_0)^{\!\top}\A^{-1}(z-x_0)
\;\ge\;\frac{q\,\dE(x_0,\Sig)^2}{2n\,\Pmax}.
\]
The constant $2n$ cannot be improved, and the bound is ratio-free: positive
definiteness of $\A$ is used only qualitatively; the constant involves $\Pmax$
alone. (Convex $\Om$ satisfies \eqref{eq:star} for every $x_0\in\Om$.)
\end{theorem}

\begin{proof}
Set $w=u+\tfrac{q}{2n}f$ with $f$ as in Proposition~\ref{prop:f}. For
$\varphi\in V$, Gauss--Green for the smooth field $\varphi\,\A\nabla f$ on the
Lipschitz $\Om$ gives
\[
\int_\Om\nabla\varphi\cdot\A\nabla f
=-2n\int_\Om\varphi
+\int_{\Gam}2\,(z-x_0)\cdot\nu\ \operatorname{tr}\varphi\,d\sigma ,
\]
the $\Sig$-part vanishing by the trace condition. Adding \eqref{eq:weak},
\[
B[w,\varphi]\;=\;\frac{q}{n}\int_\Gam (z-x_0)\cdot\nu\ \operatorname{tr}
\varphi\,d\sigma\;\ge\;0
\qquad\text{whenever }\varphi\in V,\ \varphi\ge0,
\]
by \eqref{eq:star}. Since $f$ is continuous and $\overline\Sig$ is compact,
\[
\inf_{z\in\Sig}f(z)\;=\;\min_{z\in\overline\Sig}f(z).
\]
Let $m=\tfrac{q}{2n}\inf_{\Sig}f$ and
$\varphi=(m-w)^+$. On $\Sig$, $\operatorname{tr}w=\tfrac{q}{2n}f\ge m$, so
$\operatorname{tr}\varphi=0$ and $\varphi\in V$, $\varphi\ge0$. Then
\[
0\;\le\;B[w,\varphi]\;=\;-\int_{\{w<m\}}\nabla w\cdot\A\nabla w\;\le\;0,
\]
so $\nabla\varphi\equiv0$, and the constant $\varphi$ has zero trace on
$\Sig$, forcing $\varphi\equiv0$: $w\ge m$ a.e. Since $\A$ and $q$ are
constant, $u$ is smooth in $\Om$; as $f(x_0)=0$,
$u(x_0)=w(x_0)\ge m$. The second inequality is
$\inf_{\Sig}f\ge\inf_{\Sig}|z-x_0|^2/\Pmax
=\dE(x_0,\Sig)^2/\Pmax$ by Proposition~\ref{prop:f}. Sharpness is
Proposition~\ref{prop:equality}.
\end{proof}

\begin{proposition}[Equality family]\label{prop:equality}
Equality holds in Theorem~\ref{thm:B} in each of the following cases.
\begin{enumerate}
\item[\emph{(i)}] \textbf{Ball.} $\Om=B_R(x_0)$, $\Sig=\partial B_R(x_0)$,
$\A=\Pmax I$: here $u=\dfrac{q\,(R^2-|z-x_0|^2)}{2n\Pmax}$ and
$u(x_0)=\dfrac{qR^2}{2n\Pmax}$.
\item[\emph{(ii)}] \textbf{Central reflection sectors.}
$\Om=B_R(x_0)\cap C$ with $C$ an open intersection of half-spaces whose
boundary hyperplanes pass through $x_0$; $\Sig$ the spherical part, $\Gam$ the
flat cuts, $\A=\Pmax I$. The function of \emph{(i)}, restricted, satisfies the
equation, vanishes on $\Sig$, and has conormal derivative
$\Pmax\nabla u\cdot\nu=-\tfrac qn(z-x_0)\cdot\nu=0$ on each cut; by uniqueness
it \emph{is} the torsion function, and $u(x_0)=qR^2/(2n\Pmax)$ with
\eqref{eq:star} holding with equality on $\Gam$.
\item[\emph{(iii)}] \textbf{Ellipsoids (anisotropic sharpness).}
$\Om=\{f<r^2\}$, $\Sig=\partial\Om$, $\Gam=\varnothing$: then
$u=\tfrac{q}{2n}(r^2-f)$ solves the problem and
$u(x_0)=\tfrac{q}{2n}r^2=\tfrac{q}{2n}\inf_{\Sig}f$, attaining the first
inequality of Theorem~\ref{thm:B}.
\end{enumerate}
\end{proposition}

\begin{proof}
Each candidate is smooth, lies in $V$, satisfies
$-\nabla\!\cdot\!(\A\nabla u)=q$ by Proposition~\ref{prop:f}
(e.g.\ in (iii), $-\nabla\!\cdot\!(\A\nabla u)=\tfrac{q}{2n}\cdot 2n=q$),
vanishes on $\Sig$, and has zero conormal derivative on $\Gam$ where present;
uniqueness of the weak solution identifies it with the torsion function. The
stated values are direct evaluations.
\end{proof}

\begin{corollary}[Distance to the full boundary always works]\label{cor:fullbdry}
Let $u$ be as in Definition~\ref{def:weak} with \emph{constant}
$\A\le\Pmax I$, and let the $\Sig/\Gam$ split be arbitrary. Then, for every
$x_0\in\Om$,
\[
u(x_0)\;\ge\;\frac{q\,\mathrm{dist}(x_0,\partial\Om)^2}{2n\,\Pmax},
\]
sharp on the ball with $\Sig=\partial\Om$.
\end{corollary}

\begin{proof}
Let $r=\mathrm{dist}(x_0,\partial\Om)$ and $B=B_r(x_0)\subseteq\Om$. By
Lemma~\ref{lem:compare}, $u\ge w_B$ a.e.\ on $B$, where $w_B$ is the
full-Dirichlet torsion of $B$ for $\A$. Theorem~\ref{thm:B} applied to
$(B,\Sig=\partial B,\Gam=\varnothing)$ --- the star condition is vacuous ---
gives $w_B(x_0)\ge \tfrac{q}{2n}\inf_{\partial B}f\ge qr^2/(2n\Pmax)$. Both
$u$ and $w_B$ are continuous in $B$ (constant coefficients), so the a.e.\
inequality holds at $x_0$.
\end{proof}

\begin{remark}[Metric scope]\label{rem:scope}
Theorem~\ref{thm:B} is \emph{chordal}: it uses straight-line $\A$-distance to
$\Sig$, which dominates $\dE^2/\Pmax$. The star condition \eqref{eq:star} does
not force $\dP(x_0,\Sig)=\dE(x_0,\Sig)$, so no path-metric strengthening is
asserted; likewise Corollary~\ref{cor:fullbdry} is asserted for constant
tensors only --- for variable coefficients, only the mean bound of
Theorem~\ref{thm:P1} is claimed. The contrast established by Sections
\ref{sec:horn}--\ref{sec:dichotomy}: distance to $\partial\Om$ bounds $u$
unconditionally (Corollary~\ref{cor:fullbdry}); distance to $\Sig$ does so
under the star condition --- a sharp \emph{sufficient} hypothesis whose
removal permits arbitrarily severe failure.
\end{remark}

\section{Tapered horns: closure of the distance conjecture}\label{sec:horn}

\begin{definition}[Truncated horn]\label{def:horn}
Fix $h>0$, $k>0$, $\eps\in(0,1)$, $\delta\in(0,1)$, and set
$g(x)=\eps h\,(x/h)^{k/2}$ on $[0,h]$. Define
\[
\Om_{k,\eps,\delta}\;=\;\bigl\{(x,y)\in(\delta h,h)\times\R^{n-1}\;:\;
|y|<g(x)\bigr\},\qquad
\Sig\;=\;\{h\}\times\{|y|<\eps h\},
\]
with $\Gam$ the remaining boundary: the lateral wall $\{|y|=g(x)\}$ and the end
face $F_\delta=\{\delta h\}\times\{|y|<g(\delta h)\}$. Take $\A=P\,I$ with
constant $P>0$. Since $g\in C^\infty([\delta h,h])$ has bounded slope and the
wall meets the two flat faces transversally, $\Om_{k,\eps,\delta}$ is a bounded
Lipschitz domain, so Definition~\ref{def:weak} applies. Write $\rho=|y|$.
\end{definition}

\begin{proposition}[Exact solution with exactly vanishing wall flux]\label{prop:exact}
Let
\[
a\;=\;\frac{q}{P\,\bigl(2+k(n-1)\bigr)},\qquad b\;=\;\frac{ak}{2},\qquad
v(x,y)\;=\;a\,(h^2-x^2)\;-\;b\,|y|^2 .
\]
Then, identically:
\begin{enumerate}
\item[\emph{(i)}] $-P\Delta v=q$ on $\R^n$;
\item[\emph{(ii)}] $\partial_\nu v=0$ at every point of the lateral wall
$\{\rho=g(x)\}$, $0<x<h$;
\item[\emph{(iii)}] on the end face $F_\delta$ (outward normal $-e_1$),
$\partial_\nu v=-\partial_x v=2a\delta h>0$;
\item[\emph{(iv)}] $\sup_{\overline{\Om}_{k,\eps,\delta}} v=v(\delta h,0)
= a h^2(1-\delta^2)$, and $v|_{\Sig}=-b\rho^2\in[-b(\eps h)^2,0]$.
\end{enumerate}
\end{proposition}

\begin{proof}
(i) $\Delta v=\partial_x^2 v+\Delta_y v=-2a-2b(n-1)$, so
$-P\Delta v=Pa\bigl(2+k(n-1)\bigr)=q$.
(ii) The wall is the zero set of $F(x,y)=|y|-g(x)$, with outward normal
proportional to $(-g'(x),\,y/|y|)$. Since
$\nabla v=(-2ax,\,-2by)$,
\[
\nabla v\cdot(-g',\,y/|y|)\;=\;2axg'(x)-2b\rho
\;\overset{\rho=g}{=}\;2axg'(x)-2bg(x)\;=\;\bigl(ak-2b\bigr)g(x)\;=\;0,
\]
using $xg'(x)=\tfrac{k}{2}g(x)$ for the power profile and $b=ak/2$.
(iii)--(iv) are direct evaluations; $v$ is strictly decreasing in $x$ and in
$\rho$.
\end{proof}

\begin{theorem}[Comparison]\label{thm:comparison}
On $\Om_{k,\eps,\delta}$, the torsion function satisfies
\[
0\;\le\;u\;\le\;v+b(\eps h)^2\quad\text{a.e.},
\qquad\text{hence}\qquad
\norminf{u}\;\le\;a h^2\bigl(1-\delta^2\bigr)+b(\eps h)^2 .
\]
\end{theorem}

\begin{proof}
For $\varphi\in H^1(\Om)$ the Gauss--Green identity on the Lipschitz domain,
applied to the field $\varphi\,P\nabla v$ with $v$ smooth, gives
\[
B[v,\varphi]\;=\;\int_\Om q\,\varphi
\;+\;\int_{\partial\Om}P\,(\partial_\nu v)\,\operatorname{tr}\varphi\,d\sigma .
\]
By Proposition~\ref{prop:exact}, the boundary integrand vanishes on the wall,
equals $2a\delta h\,P\operatorname{tr}\varphi$ on $F_\delta$, and for
$\varphi\in V$ the $\Sig$-part vanishes because $\operatorname{tr}\varphi=0$
there. Hence, for $\varphi\in V$ with $\varphi\ge 0$ (whence
$\operatorname{tr}\varphi\ge 0$),
\[
B[u-v,\varphi]\;=\;-\,2a\delta h\,P\int_{F_\delta}\operatorname{tr}\varphi\,
d\sigma\;\le\;0 .
\]
Set $M=b(\eps h)^2$ and $\varphi=(u-v-M)^+$. Its trace on $\Sig$ equals
$(b\rho^2-M)^+=0$ since $\rho\le\eps h$ on $\Sig$, so $\varphi\in V$, and
$\varphi\ge0$. By Stampacchia,
\[
0\;\ge\;B[u-v,\varphi]\;=\;\int_{\{u-v>M\}}\nabla(u-v)\cdot\A\nabla(u-v)
\;\ge\;0,
\]
so $\nabla\varphi\equiv 0$; the constant $\varphi$ has zero trace on $\Sig$,
hence $\varphi\equiv0$ and $u\le v+M$. Combining with Lemma~\ref{lem:pos} and
Proposition~\ref{prop:exact}(iv) yields the $L^\infty$ bound.
\end{proof}

\begin{lemma}[Distances on the horn]\label{lem:dist}
For $z=(x,y)\in\Om_{k,\eps,\delta}$,
\[
\dE(z,\Sig)\;=\;\dP(z,\Sig)\;=\;h-x,
\qquad\text{and}\qquad
\dA(z,\Sig)=\frac{h-x}{\sqrt P}.
\]
Consequently $R_{\mathrm E}=R_\Om=\sqrt{P}\,R_\A=h(1-\delta)$.
\end{lemma}

\begin{proof}
The horizontal segment $t\mapsto(t,y)$, $t\in[x,h]$, stays in $\Om$ because $g$
is nondecreasing, and ends on $\overline\Sig$; hence $\dP\le h-x$. Any path from
$z$ to $\Sig$ has length at least its $x$-displacement $h-x$; and
$\dE\ge h-x$ for the same reason while $\dE\le\dP$. The last identity is
Lemma~\ref{lem:metric}. The supremum of $h-x$ over $\Om$ is $h(1-\delta)$.
\end{proof}

\begin{theorem}[Falsification on Lipschitz domains]\label{thm:A}
On $\Om_{k,\eps,\delta}$,
\begin{equation}\label{eq:beta}
\frac{\Pmax\,\norminf{u}}{q\,R^2}
\;\le\;\beta(k,\eps,\delta,n)
\;:=\;\frac{1-\delta^2+\tfrac{k}{2}\eps^2}{\bigl(2+k(n-1)\bigr)(1-\delta)^2},
\end{equation}
where $R$ denotes any of $R_{\mathrm E}$, $R_\Om$, $\sqrt{\Pmax}R_\A$
(all equal here, with $\Pmax=P$). Since
$\beta(k,\eps,\delta,n)\to \frac{1}{2+k(n-1)}<\frac{1}{2n}$ as
$(\eps,\delta)\to(0,0)$ whenever $k>2$, every $k>2$ admits explicit
$(\eps,\delta)$ for which \emph{all three} inequalities in \eqref{eq:P2} fail.
\end{theorem}

\begin{proof}
Divide the bound of Theorem~\ref{thm:comparison} by
$qR^2/P=qh^2(1-\delta)^2/P$ and substitute $a,b$:
\[
\frac{P\norminf{u}}{qR^2}
\le \frac{P\bigl(ah^2(1-\delta^2)+b(\eps h)^2\bigr)}{qh^2(1-\delta)^2}
=\frac{(1-\delta^2)+\tfrac k2\eps^2}{(2+k(n-1))(1-\delta)^2}.
\]
For $k>2$, $\tfrac1{2+k(n-1)}<\tfrac1{2n}\iff 2n<2+k(n-1)\iff k>2$; continuity
of $\beta$ finishes. By Lemma~\ref{lem:dist} the three statements
\eqref{eq:P2} coincide on this family, so all fail together.
\end{proof}

\begin{corollary}[Exact rational instance]\label{cor:instance}
For $n=3$, $k=4$, $\delta=\tfrac{3}{20}$, $\eps=\tfrac15$:
\[
\frac{\Pmax\norminf{u}}{qR^2}
\;\le\;\frac{1-\tfrac{9}{400}+\tfrac{2}{25}}{10\cdot\bigl(\tfrac{17}{20}\bigr)^{2}}
\;=\;\frac{423/400}{2890/400}
\;=\;\frac{423}{2890}\;<\;\frac16,
\qquad \frac16-\frac{423}{2890}=\frac{88}{4335}.
\]
A bounded Lipschitz domain, constant scalar coefficient, flat Dirichlet mouth:
every hypothesis of the conjecture holds and the coefficient is $12.2\%$ below
the conjectured sharp constant.
\end{corollary}

\begin{corollary}[The infimum is zero]\label{cor:inf}
$\displaystyle
\inf\Bigl\{\Pmax\norminf{u}/(q R^2)\Bigr\}=0$, the infimum over all admissible
$(\Om,\Sig,\A,q)$, in each of the three metrics; it is not attained.
\end{corollary}

\begin{proof}
Take $\eps_k^2=1/k$ and $\delta_k=1/k$ in \eqref{eq:beta}:
$\beta=\bigl(\tfrac32-k^{-2}\bigr)/\bigl((2+k(n-1))(1-k^{-1})^2\bigr)
\to 0$ as $k\to\infty$. The ratio is strictly positive for every fixed datum by
Lemma~\ref{lem:pos}.
\end{proof}

\begin{proposition}[Cusp sharpening: the untruncated horn]\label{prop:cusp}
Let $k>2$ and let $\Om_{k,\eps}$ be the full cusp horn ($\delta=0$), with
$\Sig$ the mouth. Define
$V_0$ as the $H^1$-closure of Lipschitz functions on $\overline\Om$ vanishing
on a neighborhood of $\overline\Sig$, and let $u\in V_0$ solve \eqref{eq:weak}
for $\varphi\in V_0$. Then $u$ exists and is unique,
\[
\norminf{u}\;\le\;a h^2\Bigl(1+\tfrac{k}{2}\eps^2\Bigr),
\qquad
\frac{\Pmax\norminf{u}}{qR^2}\;\le\;\frac{1+\tfrac k2\eps^2}{2+k(n-1)},
\]
with $R=h$, and this is $<\tfrac{1}{2n}$ exactly when
$\eps^2<\dfrac{(n-1)(k-2)}{nk}$. (For $n=3$, $k=4$: $\eps^2<\tfrac13$.)
\end{proposition}

\begin{proof}
\emph{Well-posedness.} For Lipschitz $\varphi$ vanishing near $\overline\Sig$,
the segment from $(x,y)$ to $(h,y)$ lies in $\overline\Om$ ($g$ nondecreasing),
so $\varphi(x,y)=-\int_x^h\partial_t\varphi(t,y)\,dt$ and, by Cauchy--Schwarz
and Fubini, $\int_\Om\varphi^2\le h^2\int_\Om|\partial_x\varphi|^2$. The
Poincar\'e inequality passes to the closure $V_0$; Lax--Milgram applies, and
Lemma~\ref{lem:pos} holds verbatim (the truncations used there preserve the
approximating class).

\emph{Variational identity for $v$.} Fix Lipschitz $\varphi$ vanishing near
$\overline\Sig$ and $s\in(0,h)$. On the Lipschitz domain
$\Om\cap\{x>s\}$, Gauss--Green for $\varphi P\nabla v$ gives
$\int_{\{x>s\}}\nabla\varphi\cdot P\nabla v
=\int_{\{x>s\}}q\varphi+\int_{\partial}\!P(\partial_\nu v)\varphi$;
the wall term vanishes by Proposition~\ref{prop:exact}(ii), the $\Sig$ term
because $\varphi\equiv0$ there, and the face $\{x=s\}$ contributes at most
$P\sup|\nabla v|\,\|\varphi\|_\infty\,\kappa_{n-1}g(s)^{n-1}\to0$ as
$s\downarrow0$. Dominated convergence gives $B[v,\varphi]=\int q\varphi$, and by
density the identity holds for all $\varphi\in V_0$; hence $B[u-v,\varphi]=0$ on
$V_0$.

\emph{Admissible test function.} Let $M=b(\eps h)^2$ and $\psi=v+M$. On
$\overline\Om$ one has $\rho\le g(x)\le\eps h$, so
$\psi\ge a(h^2-x^2)\ge0$ everywhere. If $u_j$ are Lipschitz approximants of $u$
vanishing near $\overline\Sig$, then $(u_j-\psi)^+$ vanishes near
$\overline\Sig$ as well (there $u_j=0$ and $\psi\ge0$) and is Lipschitz; by
Stampacchia truncation continuity, $(u-\psi)^+\in V_0$.

\emph{Conclusion.} Testing $B[u-v,\cdot]=0$ with $(u-v-M)^+=(u-\psi)^+$ gives,
as in Theorem~\ref{thm:comparison}, $u\le v+M\le ah^2+b(\eps h)^2$ together
with $u\ge0$. Lemma~\ref{lem:dist} holds verbatim with $\delta=0$, so $R=h$.
Finally
\[
\frac{1+\tfrac k2\eps^2}{2+k(n-1)}<\frac1{2n}
\iff 2n+nk\eps^2<2+k(n-1)
\iff \eps^2<\frac{(n-1)(k-2)}{nk}. \qedhere
\]
\end{proof}

\begin{remark}[Status of the cusp realization]\label{rem:cuspstatus}
Proposition~\ref{prop:cusp} is logically auxiliary: every closure statement
above --- Theorem~\ref{thm:A}, Corollaries~\ref{cor:instance} and
\ref{cor:inf} --- rests entirely on the truncated Lipschitz family. The cusp
version uses the nonstandard boundary realization $V_0$ at the tip and is
flagged for independent functional-analytic review.
\end{remark}

\section{The dichotomy within the horn family: the cone as exact
transition}\label{sec:dichotomy}

\begin{proposition}[Star deficit of the horn wall]\label{prop:deficit}
On the wall of $\Om_{k,\eps,\delta}$ (or of the untruncated horn), for
$x_0=(x_*,0)$ with $\delta h\le x_*<h$, the star integrand satisfies
\[
(z-x_0)\cdot\nu\;\propto\;g(x)-\,(x-x_*)\,g'(x)
\;=\;\Bigl(1-\tfrac k2\Bigr)g(x)\;+\;x_*\,g'(x),
\]
and on the end face $F_\delta$ (normal $-e_1$),
$(z-x_0)\cdot\nu=x_*-\delta h\ge0$. Hence:
$k\le 2$ implies \eqref{eq:star} at every axis point $x_0$, while for every
$k>2$ the wall integrand is negative for all
$x>x_*\,k/(k-2)$, so \eqref{eq:star} fails at every interior $x_0$ once the
horn is long enough.
\end{proposition}

\begin{proof}
With $z=(x,g(x)\hat y)$ and outward normal $\propto(-g'(x),\hat y)$,
$(z-x_0)\cdot\nu\propto -(x-x_*)g'(x)+g(x)$; substituting
$xg'=\tfrac k2 g$ gives the display. For $k\le2$ both terms are nonnegative
($g,g'\ge0$). For $k>2$,
$(1-\tfrac k2)g+x_*g'=g\bigl(1-\tfrac k2+\tfrac{k x_*}{2x}\bigr)<0$ iff
$x>x_*k/(k-2)$.
\end{proof}

\begin{theorem}[Dichotomy and pinch at the cone]\label{thm:dichotomy}
Fix $n\ge2$ and the horn family of Definition~\ref{def:horn} with $P=\Pmax$.
\begin{enumerate}
\item[\emph{(i)}] \textbf{Sub-conical ($k\le2$): the distance bound holds with
the conjectured constant.} For every $\eps,\delta$,
\[
\norminf{u}\;\ge\;\frac{q\,R^2}{2n\,\Pmax},\qquad R=h(1-\delta).
\]
\item[\emph{(ii)}] \textbf{Super-conical ($k>2$): it fails.} For
$(\eps,\delta)$ small (explicitly, whenever
$\beta(k,\eps,\delta,n)<\tfrac1{2n}$, e.g.\ the window of
Proposition~\ref{prop:cusp}), the coefficient is strictly below
$\tfrac1{2n}$.
\item[\emph{(iii)}] \textbf{Cone ($k=2$): pinched at the conjectured constant.}
\[
\frac1{2n}\;\le\;\frac{\Pmax\norminf{u}}{qR^2}
\;\le\;\frac{1-\delta^2+\eps^2}{2n\,(1-\delta)^2}
\;\xrightarrow[(\eps,\delta)\to0]{}\;\frac1{2n}.
\]
\end{enumerate}
Within the power-law horn family, $k=2$ is thus the exact transition: the
star condition --- a sharp \emph{sufficient} hypothesis, by
Theorem~\ref{thm:B} and Proposition~\ref{prop:equality} --- controls the
sub-conical regime at the constant $1/(2n)$ in every dimension, and beyond the
cone its removal is accompanied by arbitrarily severe failure of the distance
bound (Corollary~\ref{cor:inf}). Necessity of the star condition across all
geometries is not asserted.
\end{theorem}

\begin{proof}
(i) By Proposition~\ref{prop:deficit}, every axis point
$x_0=(x_*,0)$, $x_*\in(\delta h,h)$, satisfies \eqref{eq:star};
Theorem~\ref{thm:B} with $\A=\Pmax I$ and
$\dE(x_0,\Sig)=h-x_*$ (Lemma~\ref{lem:dist}) gives
$u(x_0)\ge q(h-x_*)^2/(2n\Pmax)$. Since $u$ is continuous in $\Om$,
$\norminf{u}\ge\sup_{x_*\downarrow\delta h}q(h-x_*)^2/(2n\Pmax)
=qR^2/(2n\Pmax)$.
(ii) is Theorem~\ref{thm:A} / Proposition~\ref{prop:cusp}.
(iii) The lower bound is (i) at $k=2$; the upper bound is \eqref{eq:beta} at
$k=2$, where $2+k(n-1)=2n$.
\end{proof}

\section{Mechanism and internal consistency}\label{sec:mechanism}

\begin{remark}[Two exact one-dimensional balances; the graded-resistance
identity]\label{rem:mechanism}
Theorem~\ref{thm:oned} is the conductivity-graded balance: sink $q$ per unit
\emph{length}, conductivity $A(x)$ varying; its sharp constant is
$1/2=1/(2n)|_{n=1}$. The horns realize the volume-graded balance: sink $q$ per
unit \emph{volume} in a tube of cross-sectional volume
$\mathfrak A(s)=\kappa_{n-1}g(s)^{n-1}$. The far-end value of the exact
solution reproduces the series (graded-tube) functional identically: with
$V(s)=\int_0^s\mathfrak A$,
\[
\frac{q}{P}\int_0^h\frac{V(s)}{\mathfrak A(s)}\,ds
=\frac{q}{P}\int_0^h\frac{s}{\tfrac{k(n-1)}{2}+1}\,ds
=\frac{q\,h^2}{P\,\bigl(2+k(n-1)\bigr)}
=v(0,0).
\]
The conjectured functional $q\,d(\cdot,\Sig)^2/(2n\Pmax)$ extrapolated the
first balance; it retains the distance but discards the volume grading
$V/\mathfrak A$, and the horns are exactly the geometries on which the two
functionals separate without bound.
\end{remark}

\begin{remark}[Theorem~\ref{thm:P1} on the horns]\label{rem:p1horn}
Every ball inside a horn lies in the slab $\{|y|<\eps h\}$, so its radius is at
most $\eps h$ and Theorem~\ref{thm:P1} asserts only
$\norminf u\ge q\eps^2h^2/(n(n+2)\Pmax)$ --- far below the actual value
$\approx ah^2$ for small $\eps$, and consistent with
Theorem~\ref{thm:comparison}. The mean bound and the falsification never
touch.
\end{remark}

\begin{remark}[Corollary~\ref{cor:fullbdry} on the horns]\label{rem:fullbdryhorn}
On a horn, $\mathrm{dist}(z,\partial\Om)\le g(x)\le\eps h$ for every $z$, so
the unconditional bound of Corollary~\ref{cor:fullbdry} makes assertions only
at the scale $\eps h$, again consistent with Theorem~\ref{thm:comparison}. The
closure of the distance conjecture requires the wall to \emph{reflect}: if the
wall is made Dirichlet it joins $\Sig$, all distances to the supply collapse to
$O(\eps h)$, and no conflict can arise.
\end{remark}

\begin{remark}[Verification]\label{rem:verify}
Every displayed algebraic identity and rational inequality above --- the 1D
flux law and its bounds (Theorem~\ref{thm:oned}), the ball-mean constant
$2r^2/(n+2)$ and the resulting $n(n+2)$ (Theorem~\ref{thm:P1}; $15$ at $n=3$),
the identities of Proposition~\ref{prop:f}, the horn solution and wall-flux
identities (Proposition~\ref{prop:exact}), the coefficient formulas and
windows, the instance $423/2890<1/6$, the star deficit
(Proposition~\ref{prop:deficit}), and the cone marginality
$1/(2+k(n-1))|_{k=2}=1/(2n)$ --- has been machine-verified in exact
symbolic/rational arithmetic ($28$ independent checks; companion artifact
\texttt{p2-softcheck.py}, exit $0$). Finite-volume corroboration of the horn
coefficients exists in the companion harness and is deliberately excluded
here: this paper contains forced material only.
\end{remark}

\end{document}
